\chapter{Implementation}
\label{ch:implementation}

This chapter details the technical implementation of the MEDF platform, translating the architectural and algorithmic design into a concrete software artifact. The implementation is based on a modern Python technology stack and emphasizes code quality, testability, and reproducibility. The complete source code is available in the accompanying GitHub repository.

\section{Technology Stack}
\label{sec:tech_stack}

The project uses a curated set of open-source libraries and frameworks:

\begin{itemize}[leftmargin=2em]
  \item \textbf{Backend Framework.} FastAPI is used for the API server, providing a high-performance, asynchronous foundation with excellent data validation capabilities through its integration with Pydantic.

  \item \textbf{Frontend Framework.} Streamlit is used for the interactive web dashboard, enabling rapid development of data-centric user interfaces directly in Python.

  \item \textbf{Data Persistence.} SQLAlchemy serves as the Object-Relational Mapper (ORM) for interacting with a simple SQLite database. This choice prioritizes ease of setup and reproducibility for a research prototype, while the ORM provides a clear path for future migration to a more robust database.

  \item \textbf{Numerical and Scientific Computing.} NumPy is used for all numerical operations, particularly in the scoring and conflict detection engines, providing efficient and reliable array computations. The \texttt{pymoo} library is used for the implementation of the NSGA-II algorithm.

  \item \textbf{Configuration.} Framework definitions are managed using YAML, a human-readable data serialization format, which allows for easy editing and versioning of governance criteria.

  \item \textbf{Testing.} Pytest is used as the testing framework, with a comprehensive suite of over 50 tests covering unit, integration, property-based, and end-to-end scenarios.
\end{itemize}

\section{Codebase Structure and Critical Modules}
\label{sec:codebase}

The repository is organized into several distinct directories, reflecting the modular architecture. The core application logic resides in the \texttt{app/} directory.

\subsection{\texttt{app/main.py}}

This is the main entry point for the FastAPI application. It is responsible for initializing the application, setting up the database connection, loading the framework definitions from the YAML registry at startup, and including the API routers.

\subsection{\texttt{app/models.py}}

This module contains all the Pydantic data models that define the API schemas for requests and responses. It also includes the SQLAlchemy ORM models for the database tables (e.g., \texttt{DBStakeholderProfile}). This strict, centralized schema definition is critical for ensuring data consistency and providing clear API contracts.

\subsection{\texttt{app/scoring\_engine.py}}

This module contains the core algorithmic implementation for the scoring process. It includes functions for Likert score normalization, weight validation, and the implementations of the TOPSIS and WSM scoring methods. The code is written using NumPy for efficient vector and matrix operations.

\subsection{\texttt{app/routers/}}

The \texttt{app/routers/} directory contains the individual FastAPI routers for each major API endpoint:

\begin{itemize}[leftmargin=2em]
  \item \texttt{evaluate.py} handles the \texttt{/api/evaluate} endpoint. It receives the \texttt{EvaluateRequest} payload, computes effective weights through section-based product pooling of stakeholder and framework weights, coordinates calls to the scoring engine, and returns a validated \texttt{EvaluationResult} response.

  \item \texttt{conflicts.py} handles the \texttt{/api/conflicts} endpoint and is the \emph{de facto} conflict analysis engine. It implements the complete Spearman rank correlation computation for both weights-only and contribution-based conflict metrics, identifies divergent dimensions, determines conflict severity levels, and integrates the harm assessment module. As noted in Chapter~\ref{ch:architecture}, this logic currently remains in the router layer rather than an extracted standalone engine module.

  \item \texttt{pareto.py} handles the \texttt{/api/pareto} endpoint and implements the NSGA-II multi-objective optimization for computing Pareto-optimal consensus weight vectors.
\end{itemize}

This modular approach keeps the endpoint logic clean and organized, though the concentration of conflict analysis logic within the router rather than a dedicated engine module is an acknowledged technical debt item.

\section{Design Decisions and Implementation Rationale}
\label{sec:decisions}

A critical aspect of the implementation process was making deliberate engineering decisions to balance rigor, reproducibility, and development velocity. These decisions are documented in \texttt{docs/design\_decisions.md}.

\begin{itemize}[leftmargin=2em]
  \item \textbf{Why Weighted Scoring?} A simple unweighted average of dimension scores would fail to capture the differential importance of ethical criteria. Weighted scoring, where weights are explicitly defined by stakeholders and frameworks, provides a more nuanced and realistic evaluation model.

  \item \textbf{Why Configurable Stakeholder Weights?} Ethical priorities are subjective and context-dependent. Hardcoding a single set of weights would impose a single value system. Making weights configurable is the core mechanism that allows MEDF to model multi-stakeholder perspectives and surface conflicts.

  \item \textbf{Why a Sequential Evaluation Pipeline?} The evaluation pipeline (weighting, scoring, aggregation) is designed as a deterministic, sequential process. This ensures that the results are reproducible and the logic is easy to follow and audit. Parallelizing these steps would introduce unnecessary complexity without a clear benefit for this application.

  \item \textbf{Why a Three-Tier Architecture?} The separation of the presentation (Streamlit), API (FastAPI), and data layers creates a loosely coupled system. This allows each component to be developed, tested, and scaled independently. For example, the Streamlit frontend could be replaced with a React application without requiring any changes to the backend API.
\end{itemize}

\section{Reproducibility and Testing}
\label{sec:testing}

Ensuring the reproducibility of the research artifact was a primary implementation goal. This was achieved through several mechanisms:

\begin{itemize}[leftmargin=2em]
  \item \textbf{Dependency Pinning.} All Python dependencies are pinned to specific versions in \texttt{requirements.txt} to ensure a consistent runtime environment.

  \item \textbf{Deterministic Algorithms.} The core scoring algorithms are deterministic. The NSGA-II implementation uses a fixed random seed during evaluation runs to ensure that the Pareto frontier results are also reproducible.

  \item \textbf{Comprehensive Test Suite.} The Pytest suite includes tests that lock down the API contract (\texttt{test\_api\_contract\_lock.py}), verify the mathematical correctness of the scoring engine against hand-calculated examples (\texttt{test\_topsis\_hand\_verification.py}), and run end-to-end evaluations with the case study data to prevent regressions.

  \item \textbf{Continuous Integration.} GitHub Actions are configured to run the complete test suite on every push and pull request, ensuring that the codebase remains in a valid and tested state at all times.
\end{itemize}

This rigorous approach to implementation and testing ensures that MEDF is not just a conceptual framework, but a reliable and defensible software tool.
